Software Design Document


Revision History

Group --- Snapshot 1 Update (Month 1 - 3) ---- Version 0.1

- Added Introduction and design considerations

Group --- Snapshot 2 Update (Month 4 - 6) ---- Version 0.2

- Added Architectural Strategies, Policies and Tactics

Group --- Snapshot 3 Update (Month 7 - 9) ---- Version 0.3

- Added Detailed System Design

Group --- Snapshot 4 Update (Month 10 - 12) ---- Version 0.4

- Edited and finalized SDD

Introduction:

This document will go over the software requirements for the RoboSub project, and focuses on the
technical side. It will be intended for readers that are developers, managers, and technical 
users and is designed to help understand the code, update procedures, and architecture. The
software used will be capable of onboard computing and control, using micocontroller drives and
customized thrusters boards. 

Design Considerations:

The design of this system depends on several assumptions and constraints. The system assumes that 
the main computing device, control board, sensors, actuators, and development tools are available 
and working properly. If any of these parts change, the software design may need to be updated to 
remain compatible with the new setup. The system also has limitations that must be considered during 
development. These include limited processing power, memory, storage, testing time, and access to 
physical hardware. Because the vehicle may operate in a real-world environment, the software should 
be reliable, easy to understand, and able to handle possible errors. The design should also support 
communication between the major subsystems so that vision, control, sensors, and autonomous behavior 
can work together.

Architectural Strategies:

The system is designed using a modular architecture so that each major part of the project can be 
developed, tested, and updated separately. The main subsystems include vision, control, sensing, 
communication, simulation, and user documentation. Each subsystem should have a clear responsibility 
and should exchange information through defined interfaces.This design was chosen because it makes 
the project easier to maintain and expand. If a new sensor, actuator, or feature is added later, it 
can be connected as a separate module instead of changing the entire system. The architecture also 
supports testing in a virtual environment before using real hardware, which helps reduce risk and
allows the team to find problems earlier.

Policies and Tactics:

For this project, we will use github for version control and collaborative development. We will use 
neccessary software frameworks to provide tools for robot development to provide commuincation between 
subsystems. We will include a comprehensive README where dependencies, features, and installtion 
instructions are explained. We will imploy a variety of testing methods to ensure correct functionality,
including simulation environments and real world environments. Multiple coding languages may be used 
including python and C++, and as such we will follow these languages guidelines and document our code
as need be. 


Detailed System Design:

The system is divided into several major modules that each handle a specific responsibility. The vision 
module is responsible for collecting visual input, processing images, and identifying important objects 
in the environment. This information is then shared with the control and autonomy modules so the vehicle 
can make movement decisions based on what it detects. The control module is responsible for managing the 
vehicle's movement and translating high level decisions into commands for the physical parts of the system. 
It uses information from the vision and sensor modules to adjust direction, speed, and behavior. This module 
must respond quickly and safely because delays or incorrect commands could cause the vehicle to move incorrectly.